\documentclass[
    english,
	ruledheaders=section,%Ebene bis zu der die Überschriften mit Linien abgetrennt werden, vgl. DEMO-TUDaPub
	class=report,% Basisdokumentenklasse. Wählt die Korrespondierende KOMA-Script Klasse
	thesis={type=master},% Dokumententyp Thesis (projektseminar, proseminar, bachelor, master)
	accentcolor=2c, % Akzentfarbe FG rtm
	custommargins=false,
	marginpar=false,
	BCOR=10mm, % Bindekorrektur (etwas größere mittlere Ränder)
	parskip=half-, % Absatzkennzeichnung durch Abstand vgl. KOMA-Sript
	fontsize=11pt, % Basisschriftgröße laut Corporate Design ist mit 9pt häufig zu klein
  	twoside,
  	numbers=noendperiod,
  	fleqn,
  	captions=oneline, % Einzeilige Captions zentrieren
	captions=tableheading, % Für besseren Abstand bei TabellenÜBERschriften
  	IMRAD=false
]{tudapub}


% Der folgende Block ist nur bei pdfTeX auf Versionen vor April 2018 notwendig
\usepackage{iftex}
\ifPDFTeX
\usepackage[utf8]{inputenc}%kompatibilität mit TeX Versionen vor April 2018
\fi

%%%%%%%%%%%%%%%%%%%
%Sprachanpassung & Verbesserte Trennregeln
%%%%%%%%%%%%%%%%%%%
\usepackage[main=english, ngerman]{babel}
\usepackage[autostyle]{csquotes}% Anführungszeichen vereinfacht
\usepackage{microtype}


%%%%%%%%%%%%%%%%%%%
% Literaturverzeichnis
%%%%%%%%%%%%%%%%%%%
\catcode30=12  % Workaround for current biblatex problem (https://tex.stackexchange.com/questions/564990/error-after-miktex-reinstall-text-line-contains-an-invalid-character)
\usepackage[backend=biber,style=ieee,bibencoding=utf8]{biblatex}
\addbibresource{literature.bib}

% Additional Packages
\usepackage{acro} % Acronym management
\usepackage{amsmath, amssymb}
\usepackage{booktabs} % Professional table formatting
\usepackage{siunitx}
\usepackage{tikz}
\usepackage{graphicx} % for pdf, bitmapped graphics files
%\usepackage{cleveref} % Smart cross-referencing (load late)

% Acronym definitions (acro package)
\DeclareAcronym{mpc}{
  short = MPC ,
  long  = Model Predictive Control ,
}

\begin{document}

\Metadata{
	title=Eine minimale LaTeX-Vorlage für schriftliche (Abschluss-)Arbeiten am IAT,
	author=Rudi Ratlos,
	keywords={TU Darmstadt \sep Corporate Design \sep LaTeX},
}

\title{Eine minimale \LaTeX-Vorlage für schriftliche (Abschluss-)Arbeiten am IAT}
\subtitle{\textmd{Electrical Engineering and Information Technology program}}
\author[R. Ratlos]{Rudi Ratlos}
\reviewer*[First Examiner,Supervisor]{Prof. Dr.-Ing. Rolf Findeisen \and NameGoesHere}

\department{etit}

\addTitleBox{\includegraphics[width=\linewidth]{figures/CCPS_logo}}

%\date{\today}
\submissiondate{1. April 2020}

\lowertitleback{
	Technische Universität Darmstadt\\
	Institut für Automatisierungstechnik und Mechatronik\\
	Fachgebiet Regelungstechnik und Cyber-Physische Systeme\\
	Prof. Dr.-Ing. Rolf Findeisen\\
}

\frontmatter
\maketitle

\clearpage
\setcounter{page}{1}
\section*{Task Description}

	Provided by supervisor

 \vspace{0.5cm}
\begin{tabular}{ll}
	Start date: 	& ------ \\
	End date: 		& ------ \\
	Date of the colloquium: 	& ------ \\
\end{tabular}
\clearpage

\affidavit[digital]

\section*{Abstract}
	% Abstract structure (increasing scope):
	% 1. Context (current state of the topic)
	% 2. Problem or open research points
	% 3. Solution
	% 4. Key takeaway

\section*{Zusammenfassung}
	Additional abstract in German.

\cleardoublepage

\tableofcontents
\cleardoublepage

\mainmatter

\chapter{List of Symbols and Acronyms}

\section*{Symbols}

Matrices are set bold and uppercase, vectors bold and lowercase, and scalars italic and
lowercase. A running index is italic, a naming index upright.

\begin{tabular}{lll}
    \toprule
    Symbol & Description & Unit \\
    \midrule
    \multicolumn{3}{l}{\itshape Indices and dimensions} \\
    $i, j$ & Brain region index & -- \\
    $N$ & Number of brain regions & -- \\
    $m$ & \ac{eeg} channel index & -- \\
    $M$ & Number of \ac{eeg} channels & -- \\
    $p$ & Stimulation electrode index & -- \\
    $P$ & Number of stimulation electrodes & -- \\
    \midrule
    \multicolumn{3}{l}{\itshape Plant state and dynamics} \\
    $x_{i,1}, x_{i,2}, x_{i,3}$ & Pyramidal, excitatory and inhibitory \ac{psp} of region $i$ & \si{\milli\volt} \\
    $x_{i,4}, x_{i,5}, x_{i,6}$ & Time derivatives of the three \acp{psp} & \si{\milli\volt\per\second} \\
    $\mathbf{x}$ & Full plant state, $6N$ entries & -- \\
    $h_\mathrm{e}, h_\mathrm{i}$ & Excitatory and inhibitory \ac{psp} impulse response & \si{\milli\volt\per\second} \\
    $S(v)$ & Wave-to-pulse sigmoid & \si{\per\second} \\
    $A_i$ & Regional excitatory synaptic gain & \si{\milli\volt} \\
    $B$ & Inhibitory synaptic gain & \si{\milli\volt} \\
    $a, b$ & Excitatory and inhibitory rate constant & \si{\per\second} \\
    $\tau_\mathrm{e}, \tau_\mathrm{i}$ & Excitatory and inhibitory time constant & \si{\milli\second} \\
    $c_1, \dots, c_4$ & Intra-columnar synaptic contact numbers & -- \\
    $e_0$ & Half the maximum population firing rate & \si{\per\second} \\
    $v_0$ & Sigmoid midpoint potential & \si{\milli\volt} \\
    $r$ & Sigmoid steepness & \si{\per\milli\volt} \\
    $I$ & Constant excitatory background input & \si{\per\second} \\
    $\zeta_i$ & Gaussian white noise on the excitatory loop & \si{\milli\volt\per\second} \\
    $\sigma$ & Noise intensity & \si{\per\second} \\
    \midrule
    \multicolumn{3}{l}{\itshape Network} \\
    $\mathbf{W}$, $w_{ij}$ & Structural connectivity matrix and its entries & -- \\
    $K$ & Global coupling factor & -- \\
    $\ell_{ij}$ & Fibre tract length between regions $i$ and $j$ & \si{\milli\metre} \\
    $c$ & Axonal conduction speed & \si{\milli\metre\per\milli\second} \\
    $\mathbf{D}$, $d_{ij}$ & Conduction delay matrix and its entries & \si{\milli\second} \\
    $I_{\mathrm{coup},i}$ & Network coupling input to region $i$ & \si{\per\second} \\
    \midrule
    \multicolumn{3}{l}{\itshape Electrical interfaces} \\
    $y_{\mathrm{LFP},i}$, $\mathbf{y}_\mathrm{LFP}$ & Regional \ac{lfp} and its stacked vector & \si{\milli\volt} \\
    $\mathbf{y}_\mathrm{EEG}$ & Scalp \ac{eeg} channel amplitudes & -- \\
    $\mathbf{L}_\mathrm{EEG}$ & \ac{eeg} lead field, $M \times N$ & -- \\
    $\mathbf{L}_\mathrm{stim}$ & Stimulation lead field, $N \times P$ & -- \\
    $\mathbf{u}$ & Control current per stimulation electrode & \si{\milli\ampere} \\ % i want u_i, i \in P is the electrode
    $\mathbf{s}$, $s_i$ & Stimulation drive and its regional entries & \si{\milli\volt} \\
    \bottomrule
\end{tabular}

\begin{tabular}{lll}
    \toprule
    Symbol & Description & Unit \\
    \midrule
    \multicolumn{3}{l}{\itshape Model predictive control} \\
    $k$ & Sample index, and the control step index & -- \\
    $j$ & Step offset within the prediction horizon & -- \\
    $H$ & Prediction horizon length & -- \\
    $\mathbf{x}$, $n_x$ & Prediction-model state and its dimension & -- \\
    $\hat{\mathbf{x}}_k$ & State estimate at control step $k$ & -- \\
    $n_u$, $n_y$ & Input and output dimensions & -- \\
    $\mathbf{x}_{j|k}$ & Prediction $j$ steps ahead of step $k$ & -- \\
    $\mathbf{f}$ & One-step transition function of the prediction model & -- \\
    $\boldsymbol{\varphi}_j$ & Single-shooting rollout map over $j$ steps & -- \\
    $d_\mathrm{y}, d_\mathrm{u}$ & \ac{narx} output and input history lengths & -- \\
    $\ell$ & Stage cost & -- \\
    $\ell_\mathrm{x}$, $\ell_\mathrm{u}$ & State-dependent and input-dependent stage-cost terms & -- \\
    $V_\mathrm{f}$ & Terminal cost & -- \\
    $J_H$ & Finite-horizon cost & -- \\
    $\mathbf{Q}$, $\mathbf{R}$ & Weights of a quadratic stage cost & -- \\
    $w_1$, $w_2$ & Weights of the input penalties & -- \\
    $\mathbb{X}$, $\mathbb{U}$ & Admissible state and input sets & -- \\
    $\mathbb{X}_\mathrm{f}$ & Terminal set & -- \\
    $\mathbf{A}$, $\mathbf{B}$, $\mathbf{c}$ & State matrix, input matrix, and offset of an affine prediction model & -- \\
    $\mathbf{A}_q$, $\mathbf{B}_q$ & State and input Jacobians at predicted step $q$ & -- \\
    $\mathbf{I}_n$ & Identity matrix of size $n \times n$ & -- \\
    $\delta_{j|k}$ & Binary input-activation indicator & -- \\
    $u_\mathrm{max}$ & Per-electrode current bound & \si{\milli\ampere} \\
    \bottomrule
\end{tabular}

\begin{tabular}{lll}
    \toprule
    Symbol & Description & Unit \\
    \midrule
    \multicolumn{3}{l}{\itshape Short-time Fourier transform} \\
    $f_\mathrm{s}$ & Sampling rate of the segmented signal & \si{\hertz} \\
    $\Delta t$, $t_k$ & Sampling interval and the time of sample $k$ & \si{\second} \\
    $n$ & Sample index within a segment & -- \\
    $q$ & Segment index on the hop grid & -- \\
    $L_\mathrm{s}$ & Segment length & -- \\
    $H_\mathrm{s}$ & Hop between consecutive segments & -- \\
    $g_n$ & Hann taper applied to a segment & -- \\
    $\nu$, $f_\nu$ & Frequency bin index and the frequency it sits at & --, \si{\hertz} \\
    $\Delta f$ & Bin spacing, $f_\mathrm{s} / L_\mathrm{s}$ & \si{\hertz} \\
    $Y_q(f_\nu)$ & \Ac{stft} coefficient of segment $q$ & -- \\
    $\hat{\Phi}_q$, $\bar{\Phi}_q$ & Periodogram and its trailing average & -- \\
    $\alpha_\nu$ & One-sided folding factor of the periodogram & -- \\
    $\kappa_{q'}$, $L_\mathrm{k}$ & Averaging kernel weights and width & -- \\
    \bottomrule
\end{tabular}

\section*{Acronyms}

\printacronyms[heading=none]

\chapter{Introduction}
% Context
% Problem Statement and Open Research Points
% Proposed Solution
% Contributions
% Thesis Structure

\chapter{Foundations}

% need introduction here

\section{Whole-Brain Biophysical Modeling}

The controllers designed in this work need to be evaluated through simulations against a ground thruth model plant.
The plant used throughout this work is a whole-brain network of Jansen-Rit neural masses coupled
through a structural connectome, following the epilepsy model of Yu et al.\ \cite{Yu2024}.
This section develops that model from the single cortical column to the network, and closes with the two electrical interfaces the controller acts through: the scalp \ac{eeg} it measures and the stimulation drive it applies.

\subsection{The Jansen-Rit Neural Mass Model}

A neural mass model abandons the individual neuron and describes a cortical macro-column
by the mean membrane potential and mean firing rate of whole populations, an approximation
justified when the quantity of interest --- the extracellular field --- is itself a % no em dash
spatial average over millions of cells \cite{LopesDaSilva1974,Nunez2006}. Lopes da Silva
et al.\ introduced the formalism to explain the thalamic alpha rhythm \cite{LopesDaSilva1974},
and Jansen and Rit gave it the three-population form used here, which reproduces both
resting alpha activity and visual evoked potentials from one parameter set
\cite{Jansen1995}.
Its appeal for this work is the combination of physiological
interpretability as well as the ability to produce waveforms characteristic of epileptic seizures \cite{Wendling2002,Yu2024}.
This differentiates it from simpler models such as Wilson-Cowan or FitzHugh-Nagumo, both only two instead of six states, separate excitation from inhibition but collapse the synaptic filtering setting the frequency of the resulting
rhythm, and produce oscillations without the spike-wave morphology making a simulated
seizure recognisable \cite{David2003}.

Each node $i$ comprises a population of \acp{py} that projects to, and receives
feedback from, a population of \acp{ein} and a population of \acp{iin}, wired as shown in
Figure~\ref{fig:microcircuit}.
The excitatory branch forms a positive feedback loop and the inhibitory
branch a negative one, and the competition between the two produces the oscillations the
model is built to explain. The strengths of the four intra-columnar projections are the
average synaptic contact numbers $c_1$ to $c_4$, held at the standard proportions
$c_2 = 0.8\,c_1$ and $c_3 = c_4 = 0.25\,c_1$ so that a single number $c_1$ scales the
whole microcircuit \cite{Jansen1995}.

\begin{figure}[tb]
    \centering
    % PLACEHOLDER -- see figures/03_foundations/PLACEHOLDER_jansen_rit_microcircuit.md
    \framebox[0.8\linewidth]{\rule{0pt}{5cm}\itshape Jansen-Rit microcircuit block diagram}
    \caption{Block diagram of one Jansen-Rit node. The \ac{py} population excites the
    \acp{ein} through $c_1$ and the \acp{iin} through
    $c_3$, and receives excitatory feedback through $c_2$ and inhibitory feedback through
    $c_4$. Each population filters its incoming firing rate with the postsynaptic potential
    kernel $h_\mathrm{e}$ or $h_\mathrm{i}$ and converts the resulting potential back to a
    firing rate with the sigmoid $S$. The background input $I$, the network input
    $I_{\mathrm{coup},i}$ and the noise $\zeta_i$ enter the excitatory branch, the
    stimulation drive $s_i$ enters at the pyramidal soma, and the node output
    $y_{\mathrm{LFP},i} = x_{i,2} - x_{i,3}$ is tapped there.}
    \label{fig:microcircuit}
\end{figure}

Every population converts its input in two stages. The first is a linear pulse-to-wave
operator turning an incoming mean firing rate into an average \ac{psp} by convolution
with the impulse response
%
\begin{equation}
    h_\mathrm{e}(t) = A\,a\,t\,\mathrm{e}^{-a t}, \qquad
    h_\mathrm{i}(t) = B\,b\,t\,\mathrm{e}^{-b t}, \qquad t \geq 0,
    \label{eq:psp_kernel}
\end{equation}
%
in which the synaptic gains $A$ and $B$ set the amplitude of the excitatory and inhibitory
\ac{psp} in \si{\milli\volt} and the rate constants $a$ and $b$ set the reciprocal
lumped time constants of dendritic delay and membrane decay. The second stage is a static
wave-to-pulse operator mapping the mean membrane potential $v$ of a population onto its
mean firing rate through the sigmoid
%
\begin{equation}
    S(v) = \frac{2 e_0}{1 + \exp\!\big(r (v_0 - v)\big)},
    \label{eq:sigmoid}
\end{equation}
%
whose ceiling $2 e_0$ is the maximum population firing rate, whose midpoint $v_0$ is the
potential at which half the population fires, and whose steepness $r$ reflects the spread of
firing thresholds across the cells of the population.

Because the kernel in \eqref{eq:psp_kernel} is the Green's function of a second-order linear
operator, the convolution can be written as a pair of first-order differential equations, and
the three populations together give the six states $x_{i,1}$ to $x_{i,6}$ per node. The
potentials $x_{i,1}$, $x_{i,2}$ and $x_{i,3}$ are the \ac{psp} of the pyramidal, excitatory and
inhibitory branch, and $x_{i,4}$, $x_{i,5}$ and $x_{i,6}$ are their time derivatives, so that
the node obeys
%
\begin{align}
    \dot{x}_{i,1} &= x_{i,4}, \nonumber \\
    \dot{x}_{i,4} &= A_i a\, S\!\left(x_{i,2} - x_{i,3} + s_i\right) - 2 a\, x_{i,4} - a^2 x_{i,1}, \nonumber \\
    \dot{x}_{i,2} &= x_{i,5}, \nonumber \\
    \dot{x}_{i,5} &= A_i a \left[ I + c_2 S\!\left(c_1 x_{i,1}\right) + I_{\mathrm{coup},i} \right] - 2 a\, x_{i,5} - a^2 x_{i,2} + \zeta_i, \label{eq:jr_node} \\
    \dot{x}_{i,3} &= x_{i,6}, \nonumber \\
    \dot{x}_{i,6} &= B b\, c_4 S\!\left(c_3 x_{i,1}\right) - 2 b\, x_{i,6} - b^2 x_{i,3}, \nonumber
\end{align}
%
where $I$ is the constant excitatory background input standing in for afferents the model does
not resolve, $I_{\mathrm{coup},i}$ the input from the rest of the network developed in
Section~\ref{sec:connectome}, and $s_i$ the stimulation drive of Section~\ref{sec:electro}.
The term $\zeta_i$ is Gaussian white noise of intensity $\sigma$ entering the excitatory loop.
The noise is a necessary component. A node sitting just below its bifurcation has two
coexisting attractors, and only the noise pushes it from the quiet one into the
oscillating one, so the noise is the mechanism by which
seizures start at irregular times and by which repeated runs of one configuration differ
\cite{Yu2024}. Table~\ref{tab:jr_params} collects the parameter values, which are those of the
original single-column fit \cite{Jansen1995} except where the network model of Yu et al.\
prescribes otherwise \cite{Yu2024}.

The quantity read out of a node is the net depolarisation of the pyramidal population,
%
\begin{equation}
    y_{\mathrm{LFP},i}(t) = x_{i,2}(t) - x_{i,3}(t),
    \label{eq:lfp}
\end{equation}
%
the difference between the excitatory and inhibitory \ac{psp} arriving at the pyramidal
soma. Its physical interpretation is taken up in Section~\ref{sec:electro}; for the dynamics
it matters that $y_{\mathrm{LFP},i}$ is also the argument of the pyramidal sigmoid in
\eqref{eq:jr_node}, so the output of a node and its own feedback are the same signal.

\subsection{Regional Excitability and Seizure Onset}
\label{sec:excitability}

The excitatory synaptic gain $A_i$ is the parameter deciding which regime a node
settles into, which is why it carries a regional index while every other
parameter in \eqref{eq:jr_node} is shared. Increasing $A_i$ strengthens the excitatory
feedback loop relative to the inhibitory one, and the node passes through a Hopf
bifurcation from a stable fixed point, around which noise produces low-amplitude alpha-band
fluctuations, to a large-amplitude limit cycle of spike-wave morphology. Following the
excitability partition established for whole-brain epilepsy models
\cite{Yu2024}, three levels are used. Healthy background tissue takes
$A_i = 3.25\,\si{\milli\volt}$ and rests at the fixed point. The \ac{ez} takes
$A_i = 3.6\,\si{\milli\volt}$ and oscillates autonomously, needing no input from the
network to seize. The \ac{pz} takes the intermediate $A_i = 3.4\,\si{\milli\volt}$, which
leaves it below its own bifurcation but close enough that coupling from an active \ac{ez}
recruits it.
Figure~\ref{fig:bifurcation} shows the three levels as time series and power spectra, and contrasts a \ac{pz} node in isolation with the same node embedded in the network.

\begin{figure}[tb]
    \centering
    \includegraphics[width=0.9\linewidth]{figures/03_foundations/excitability_regimes}
    \caption{Jansen-Rit dynamics at the three excitability levels. Left, the local field
    potential $y_\mathrm{LFP}$ over \SI{4}{\second} of simulation following a
    \SI{4}{\second} transient; right, the corresponding power spectral density. Rows (a),
    (b) and (d) are uncoupled nodes: background tissue and the \ac{pz} both fluctuate
    around their fixed point at low amplitude, while the \ac{ez} has crossed into a
    large-amplitude spike-wave limit cycle with pronounced harmonics. Row (c) is the same
    \ac{pz} node embedded in the 76-region network alongside three \ac{ez} nodes, where
    coupling recruits it into an oscillation of \ac{ez}-like amplitude, the propagation
    the controller has to prevent. All runs use the parameters of
    Table~\ref{tab:jr_params} and the same noise seed.}
    \label{fig:bifurcation}
\end{figure}

\subsection{Connectome Coupling and Conduction Delays}
\label{sec:connectome}

Scaling from one column to a brain means specifying which nodes talk to which, how
strongly, and how long the signal takes to arrive. The first two are carried by the
structural connectivity matrix $\mathbf{W} \in \mathbb{R}^{N \times N}$ over $N$ regions,
whose entry $w_{ij}$ is the density of fibres between regions $i$ and $j$; the matrix is
symmetric, since diffusion-weighted tractography and the tracer collations it is validated
against recover fibre bundles without their direction \cite{Kotter2004,SanzLeon2013}. A
single global coupling factor $K$ scales the whole matrix and is the one free network
parameter, $\mathbf{W}$ itself being measured rather than fitted.

Long-range cortico-cortical projections are axonal, so what one region transmits to another
is a firing rate rather than a potential, and it arrives late. The delay follows from the
fibre length $\ell_{ij}$ between the two regions and a conduction speed $c$ assumed uniform
across the white matter, giving the delay matrix $\mathbf{D} \in \mathbb{R}^{N \times N}$
with entries $d_{ij} = \ell_{ij} / c$. Delayed presynaptic firing enters the excitatory
\ac{psp} channel of the receiving node, so the coupling term of \eqref{eq:jr_node} is
%
\begin{equation}
    I_{\mathrm{coup},i}(t) = K \sum_{j=1}^{N} w_{ij}\, S\!\big(y_{\mathrm{LFP},j}(t - d_{ij})\big).
    \label{eq:coupling}
\end{equation}
%
Together, \eqref{eq:jr_node} and \eqref{eq:coupling} constitute a stochastic delay
differential equation in the $6N$-dimensional state
$\mathbf{x}(t) = \big[x_{1,1}(t), \dots, x_{N,6}(t)\big]^\top$, whose solution from a given
initial history is the Plant every later chapter treats as ground truth.

With $\ell_{ij}$ on the order of \SI{100}{\milli\metre} and $c$ on the order of
\SI{50}{\milli\metre\per\milli\second}, inter-regional delays reach several milliseconds and
are comparable to the membrane time constants in Table~\ref{tab:jr_params}, so they shape
the phase relations deciding whether an active \ac{ez} recruits a neighbour.

\subsection{Electrophysiological Interpretation}
\label{sec:electro}

So far $y_{\mathrm{LFP},i}$ has been a state combination. Its physical reading is that
synaptic currents entering the apical dendrites of a pyramidal population, together with
the return currents at the soma, form a current dipole oriented along the somato-dendritic
axis of the cell; because pyramidal cells are aligned perpendicular to the cortical
sheet, these dipoles sum coherently across a region instead of cancelling, and they, not
the action potentials, dominate the extracellular field \cite{Nunez2006}. The net somatic depolarisation $x_{i,2} - x_{i,3}$ is the only variable of the model tracking
that current balance, and it is taken as a proxy for the regional dipole moment up to an
undetermined constant of proportionality, the convention throughout neural mass modelling
of \ac{eeg} \cite{David2003,SanzLeon2015}. The convention is a modelling choice rather
than a derivation: the model carries no cortical geometry, so it fixes neither the sign
nor the scale of the dipole, and \eqref{eq:lfp} therefore predicts the waveform of the
regional \ac{lfp} but not its absolute amplitude.

Scalp \ac{eeg} follows from the \ac{lfp} by volume conduction through the head. Quasi-static
Maxwell equations make the map linear, so with
$\mathbf{y}_\mathrm{LFP}(t) \in \mathbb{R}^{N}$ collecting the regional \acp{lfp} the $M$
channel voltages are
%
\begin{equation}
    \mathbf{y}_\mathrm{EEG}(t) = \mathbf{L}_\mathrm{EEG}\, \mathbf{y}_\mathrm{LFP}(t),
    \label{eq:eeg}
\end{equation}
%
where the \ac{eeg} lead field $\mathbf{L}_\mathrm{EEG} \in \mathbb{R}^{M \times N}$ holds in
entry $(m,i)$ the voltage electrode $m$ records per unit dipole moment in region $i$
\cite{SanzLeon2015}. The dipole scale being undetermined, $\mathbf{y}_\mathrm{EEG}$ carries
the arbitrary amplitude units of $\mathbf{y}_\mathrm{LFP}$ rather than volts, which is why
every threshold and cost in the later chapters is stated relative to the background level
of the model itself.

Stimulation runs in the opposite direction. Transcranial currents injected at $P$ scalp
electrodes produce an electric field in the head, and the component of that field along the
somato-dendritic axis polarises the pyramidal membrane, shifting its operating point on the
sigmoid without injecting a firing rate of its own. The perturbation therefore enters
\eqref{eq:jr_node} inside the pyramidal sigmoid as the additive drive $s_i$, the
macroscopic $\lambda_\mathrm{E}$ formulation of \ac{tes} adopted by large-scale modelling
studies \cite{Kunze2016,Merlet2013,Yu2024}. Collecting the electrode currents in
$\mathbf{u} \in \mathbb{R}^{P}$ and the regional drives in $\mathbf{s} \in \mathbb{R}^{N}$,
the drive is
%
\begin{equation}
    \mathbf{s}(t) = \mathbf{L}_\mathrm{stim}\, \mathbf{u}(t), \qquad
    \mathbf{1}^\top \mathbf{u}(t) = 0,
    \label{eq:stim}
\end{equation}
%
with the stimulation lead field $\mathbf{L}_\mathrm{stim} \in \mathbb{R}^{N \times P}$ and
the zero-sum constraint imposed by Kirchhoff's current law, which forbids current from accumulating in the head.

The two lead fields in \eqref{eq:eeg} and \eqref{eq:stim} describe the same physics viewed
from opposite ends. Helmholtz reciprocity states that the lead field relating a source in
the brain to a pair of scalp electrodes is identical to the field the same electrode pair
would produce in the brain when driven by unit current \cite{Malmivuo1995}, so
$\mathbf{L}_\mathrm{stim}$ and $\mathbf{L}_\mathrm{EEG}^\top$ ought to agree up to scaling.
How the two lead fields are obtained in practice is taken up in Chapter~\ref{ch:methodology}.

\begin{table}[tb]
    \centering
    \caption{Parameters of the Jansen-Rit network model. Values follow the original
    single-column fit \cite{Jansen1995} and the whole-brain epilepsy model of Yu et al.\
    \cite{Yu2024}. The global coupling factor $K$ and the noise intensity $\sigma$ deviate
    from the values of Yu et al.\ and were recalibrated for the connectome used here, as
    Chapter~\ref{ch:methodology} sets out. Rate constants are given both as inverse time
    constants and as the corresponding time constants.}
    \label{tab:jr_params}
    \begin{tabular}{llll}
        \toprule
        Symbol & Description & Value & Unit \\
        \midrule
        $A_i$   & Excitatory synaptic gain (background / \ac{pz} / \ac{ez}) & 3.25 / 3.4 / 3.6 & \si{\milli\volt} \\
        $B$     & Inhibitory synaptic gain                        & 22    & \si{\milli\volt} \\
        $a$     & Excitatory rate constant ($\tau_\mathrm{e} = \SI{10}{\milli\second}$)  & 100 & \si{\per\second} \\
        $b$     & Inhibitory rate constant ($\tau_\mathrm{i} = \SI{20}{\milli\second}$)  & 50  & \si{\per\second} \\
        $c_1$   & Pyramidal-to-excitatory synaptic contacts        & 135   & -- \\
        $c_2$   & Excitatory-to-pyramidal synaptic contacts        & 108   & -- \\
        $c_3$   & Pyramidal-to-inhibitory synaptic contacts        & 33.75 & -- \\
        $c_4$   & Inhibitory-to-pyramidal synaptic contacts        & 33.75 & -- \\
        $e_0$   & Half the maximum population firing rate          & 2.5   & \si{\per\second} \\
        $v_0$   & Sigmoid midpoint potential                       & 6     & \si{\milli\volt} \\
        $r$     & Sigmoid steepness                                & 0.56  & \si{\per\milli\volt} \\
        $I$     & Constant excitatory background input             & 90    & \si{\per\second} \\
        $\sigma$ & Noise intensity on the excitatory loop          & 280   & \si{\per\second} \\
        $K$     & Global coupling factor                           & 0.60  & -- \\
        $c$     & Axonal conduction speed                          & 50    & \si{\milli\metre\per\milli\second} \\
        \bottomrule
    \end{tabular}
\end{table}

The Plant is now fully specified, and with it the difficulty the rest of the thesis
addresses. The controller must regulate a $6N$-dimensional delayed stochastic system while
observing only $M$ linear mixtures of it and acting only through $P$ scalp currents, and it
must do so fast enough to intervene before recruitment spreads. The following section
develops the control formulation meeting those requirements, and states in doing so what it
demands of a prediction model; Section~
ef{sec:sysid} then turns to identifying such a
model from data, since solving the plant equations online is out of reach.

\section{Model Predictive Control}
\label{sec:mpc}

\Ac{mpc} chooses the electrode currents by solving, at every sampling instant, an \ac{ocp}
over a finite window of the future, applying the first input of the solution, and repeating
the computation one step later from the new measurement. Two properties suit it to this
application. The first is optimality. The controller has to suppress pathological activity
on as little delivered current as it can manage, and in neurostimulation that economy is an
objective of its own, since the delivered charge bounds both the exposure of the tissue and
the energy the device draws. \Ac{mpc} minimises that objective directly, where a fixed
feedback law spends whatever effort its gain dictates. The second is constraint handling.
Scalp stimulation is bounded by a Control Budget on per-electrode and total current, and
Kirchhoff's law of \eqref{eq:stim} forces the currents to sum to zero at every instant. \Ac{mpc} enforces both inside the optimisation, whereas a
fixed feedback law can only be clipped after the fact \cite{Rawlings2017,Mayne2000}.

\subsection{Receding Horizon Control}
\label{sec:rhc}

Let the prediction model be a discrete-time map
%
\begin{equation*}
    \mathbf{x}_{k+1} = \mathbf{f}(\mathbf{x}_k, \mathbf{u}_k),
\end{equation*}
%
with $\mathbf{x}_k \in \mathbb{R}^{n_x}$ the model state at step $k$, $\mathbf{u}_k \in
\mathbb{R}^{n_u}$ the input applied over the interval following that step, and $\mathbf{f}$
the one-step transition. Throughout this section $\mathbf{x}$ denotes the state of the
prediction model and not the plant state of Section~\ref{sec:connectome}. The two coincide
only in the hypothetical case of a controller given the plant equations; for a model
identified from data they need not share a dimension, let alone a physical meaning, and
Section~\ref{sec:output_feedback} gives a model whose state is a window of past
measurements.

At step $k$ the controller solves, over the Control Horizon of $H$ steps, the \ac{ocp}
%
\begin{align}
    \min_{\mathbf{u}_{0|k}, \dots, \mathbf{u}_{H-1|k}} \quad
    & J = \sum_{j=0}^{H-1} \ell\!\left(\mathbf{x}_{j|k}, \mathbf{u}_{j|k}\right)
      + V_\mathrm{f}\!\left(\mathbf{x}_{H|k}\right) \nonumber \\
    \text{s.t.} \quad
    & \mathbf{x}_{j+1|k} = \mathbf{f}\!\left(\mathbf{x}_{j|k}, \mathbf{u}_{j|k}\right),
      \qquad j = 0, \dots, H-1, \nonumber \\
    & \mathbf{x}_{0|k} = \hat{\mathbf{x}}_k, \label{eq:ocp} \\
    & \mathbf{g}\!\left(\mathbf{x}_{j|k}, \mathbf{u}_{j|k}\right) = \mathbf{0}, \nonumber \\
    & \mathbf{h}\!\left(\mathbf{x}_{j|k}, \mathbf{u}_{j|k}\right) \leq \mathbf{0}, \nonumber
\end{align}
%
in which the double index $j|k$ denotes the quantity predicted $j$ steps ahead from the
information available at step $k$, $\hat{\mathbf{x}}_k$ is the state the controller believes
the system to be in, $\ell$ is the stage cost, $V_\mathrm{f}$ the terminal cost, and
$\mathbf{g}$ and $\mathbf{h}$ collect the equality and inequality constraints. Simple bounds
$\mathbf{u} \in \mathbb{U}$ and $\mathbf{x} \in \mathbb{X}$ are the special case of
$\mathbf{h}$ used most often, and the zero-sum condition on the electrode currents is an
instance of $\mathbf{g}$. Writing $\mathbf{u}^\star_{0|k}, \dots, \mathbf{u}^\star_{H-1|k}$
for the minimiser, the controller discards all but the first element and applies
%
\begin{equation*}
    \mathbf{u}_k = \mathbf{u}^\star_{0|k},
\end{equation*}
%
then advances to step $k+1$, takes a fresh measurement, and solves \eqref{eq:ocp} again over
a window shifted one step into the future. Figure~\ref{fig:receding_horizon} shows the
mechanism over three consecutive steps.

\begin{figure}[tb]
    \centering
    \includegraphics[width=0.9\linewidth]{figures/03_foundations/receding_horizon}
    \caption{The receding horizon principle over three consecutive control steps. At each
    step the controller plans an input sequence over the Control Horizon of $H$ steps and
    the trajectory it predicts, applies only the first input of that sequence, and discards
    the rest. One step later the window has moved forward and the plan is recomputed from
    the new measurement, so the realised trajectory is stitched together from the first step
    of many different plans and never follows any one of them to its end. The plans depart
    from the realised trajectory because the planner uses a nominal model while the plant is
    disturbed, which is the mismatch the repetition corrects. The input is held constant
    across a sampling interval and respects the bound $\mathbb{U}$ at every step. The
    example uses a scalar system and serves only to illustrate the bookkeeping.}
    \label{fig:receding_horizon}
\end{figure}

Solving \eqref{eq:ocp} once yields an open-loop input sequence. Feedback arises only from
the repetition: because $\hat{\mathbf{x}}_k$ is refreshed from measurement before every
solve, a deviation between what the model predicted and what the plant did is corrected at
the next step rather than accumulating over the horizon. Two of the difficulties in this
work are absorbed by that mechanism. The plant of Section~\ref{sec:connectome} is driven by
the noise $\zeta_i$, while the controller optimises a single noise-free rollout of its
model. Replacing the stochastic system by the deterministic one obtained from setting every
future disturbance to its expected value is the certainty-equivalent approximation, and it
is only an approximation, because the cost of the expected trajectory is not the expected
cost of the trajectory whenever the dynamics or the cost are nonlinear. Stochastic and
robust \ac{mpc} instead propagate the disturbance distribution or a worst-case tube through
the horizon and tighten the constraints accordingly \cite{Mesbah2016}, at a cost in problem
size not taken on here. The second difficulty is model mismatch. The prediction model is
identified from data and is wrong in ways no disturbance model captures, so its rollout
degrades with $j$; the receding horizon limits the damage by acting on the near end of the
horizon, where the prediction is most trustworthy, and re-planning before the far end is
reached \cite{Hewing2020}.

Repetition alone establishes no stability guarantee. Nominal stability results for \ac{mpc}
are built on terminal ingredients, a terminal cost $V_\mathrm{f}$ chosen as a local control
Lyapunov function together with a terminal set the state is constrained to reach, which make
the optimal cost a Lyapunov function of the closed loop and give recursive feasibility
\cite{Mayne2000}. Constructing them requires an equilibrium to build the terminal set around
and a model accurate enough for that set to remain invariant under the true dynamics.
Neither is available to a data-driven controller acting on a seizure, so the closed loop of
this work carries no certificate and its stability is assessed empirically, in
closed-loop simulation. The relevant theory for that setting is the
analysis of receding horizon control without terminal ingredients, where closed-loop
performance is bounded in terms of the horizon length and a controllability condition, and a
sufficiently long horizon recovers a fraction of the infinite-horizon performance
\cite{Gruene2017}. Horizon length thus carries the burden terminal ingredients would
otherwise carry, which is why Chapter~\ref{ch:methodology} treats it as a tuning parameter
rather than fixing it by design.

\subsection{Linear and Nonlinear MPC}

The character of \eqref{eq:ocp} is decided by $\mathbf{f}$ and $\ell$. If the model is
linear, $\mathbf{x}_{k+1} = \mathbf{A}\mathbf{x}_k + \mathbf{B}\mathbf{u}_k$ with the system
matrix $\mathbf{A}$ and input matrix $\mathbf{B}$, the stage cost quadratic,
$\ell(\mathbf{x}, \mathbf{u}) = \mathbf{x}^\top \mathbf{Q} \mathbf{x} + \mathbf{u}^\top
\mathbf{R} \mathbf{u}$ with the weights $\mathbf{Q} \succeq 0$ and $\mathbf{R} \succ 0$, and
the constraints affine, then the predicted states are affine in the input sequence and can
be substituted out. What remains is a \ac{qp} in the inputs alone: convex, with a single
global minimiser a solver reaches reliably and in a number of iterations that can be bounded
in advance \cite{Rawlings2017}.

Those properties are worth striving for, and the obstacle is that the plant is not linear.
The Jansen-Rit node is nonlinear through the sigmoid of \eqref{eq:sigmoid}, and the
stimulation drive enters inside that sigmoid. Nonlinear dynamics do not by themselves force
a nonlinear controller, however, because linearity is a property of the coordinates the
model is written in as much as of the system. Koopman operator theory makes the point
precise. The flow of a nonlinear system acts linearly on functions of its state, so a model
predicting a suitably chosen set of observables forward can be linear even though the
underlying dynamics are not \cite{Rowley2009,Brunton2016}. A linear model in such a lifted
space returns \eqref{eq:ocp} to a \ac{qp}, and Section~\ref{sec:sysid} takes up the question
of whether one identified from data predicts well enough over a whole Control Horizon.

Where no such model is available, \eqref{eq:ocp} is a nonconvex \ac{nlp} and \ac{nmpc} loses
every property the linear case supplies: a solver returns a local minimiser whose quality
depends on the initial guess, several such minimisers may exist, and the number of
iterations needed to reach one is not known in advance \cite{Findeisen2002}. Choosing
between the two formulations trades the reliability of the optimisation against the fidelity
of the model, and this work implements both to see where each of them falls.

\subsection{Transcription and Numerical Solution}
\label{sec:transcription}

Turning \eqref{eq:ocp} into something a numerical solver accepts is called transcription,
and the choice is which of the predicted quantities become decision variables.

In direct single shooting the inputs alone are optimised. The dynamics are used to simulate
the trajectory forward from $\hat{\mathbf{x}}_k$, so that every predicted state is an
explicit function $\boldsymbol{\varphi}_j$ of the initial state and the inputs preceding it,
%
\begin{equation*}
    \boldsymbol{\varphi}_0 = \hat{\mathbf{x}}_k, \qquad
    \boldsymbol{\varphi}_{j+1}\!\left(\hat{\mathbf{x}}_k, \mathbf{u}_{0|k}, \dots,
    \mathbf{u}_{j|k}\right)
    = \mathbf{f}\!\left(\boldsymbol{\varphi}_j, \mathbf{u}_{j|k}\right),
\end{equation*}
%
and substituting it into \eqref{eq:ocp} leaves
%
\begin{align}
    \min_{\mathbf{u}_{0|k}, \dots, \mathbf{u}_{H-1|k}} \quad
    & \sum_{j=0}^{H-1} \ell\!\left(\boldsymbol{\varphi}_j, \mathbf{u}_{j|k}\right)
      + V_\mathrm{f}\!\left(\boldsymbol{\varphi}_H\right) \nonumber \\
    \text{s.t.} \quad
    & \mathbf{g}\!\left(\boldsymbol{\varphi}_j, \mathbf{u}_{j|k}\right) = \mathbf{0},
      \label{eq:single_shooting} \\
    & \mathbf{h}\!\left(\boldsymbol{\varphi}_j, \mathbf{u}_{j|k}\right) \leq \mathbf{0},
      \nonumber
\end{align}
%
in $H n_u$ variables. The dynamic constraints have disappeared, satisfied by construction,
so every solver iterate is already a valid trajectory of the model.

In direct multiple shooting the horizon is cut into intervals and the state at the start of
each interval joins the decision vector, with continuity imposed as an explicit equality
constraint \cite{Bock1984}. Taking every step as an interval, the problem is
%
\begin{align}
    \min_{\mathbf{x}_{0|k}, \dots, \mathbf{x}_{H|k},\; \mathbf{u}_{0|k}, \dots,
    \mathbf{u}_{H-1|k}} \quad
    & \sum_{j=0}^{H-1} \ell\!\left(\mathbf{x}_{j|k}, \mathbf{u}_{j|k}\right)
      + V_\mathrm{f}\!\left(\mathbf{x}_{H|k}\right) \nonumber \\
    \text{s.t.} \quad
    & \mathbf{x}_{j+1|k} - \mathbf{f}\!\left(\mathbf{x}_{j|k}, \mathbf{u}_{j|k}\right)
      = \mathbf{0}, \qquad
      \mathbf{x}_{0|k} - \hat{\mathbf{x}}_k = \mathbf{0}, \label{eq:multiple_shooting} \\
    & \mathbf{g}\!\left(\mathbf{x}_{j|k}, \mathbf{u}_{j|k}\right) = \mathbf{0}, \qquad
      \mathbf{h}\!\left(\mathbf{x}_{j|k}, \mathbf{u}_{j|k}\right) \leq \mathbf{0}, \nonumber
\end{align}
%
in $H n_u + (H+1) n_x$ variables. The continuity equalities hold only at convergence, so an
intermediate iterate consists of trajectory pieces separated by gaps at the interval
boundaries, as Figure~\ref{fig:shooting} contrasts with the single unbroken rollout of
\eqref{eq:single_shooting}.

The path constraints $\mathbf{g}$ and $\mathbf{h}$ appear in both transcriptions in different
forms. Multiple shooting imposes them on decision variables directly, so a bound on a
predicted state is a bound on a variable and the Jacobian of the constraint set is block
banded, each block coupling one interval to the next. Under single shooting the states are
not variables. A constraint on $\boldsymbol{\varphi}_j$ is a nonlinear function of the $j$
inputs preceding it, its Jacobian row fills in as the horizon advances, and a simple state
bound late in the horizon ends up a dense nonlinear inequality. Constraints on the input
escape this. Kirchhoff's law of \eqref{eq:stim} and the Control Budget involve no predicted
state, so single shooting leaves them the simple bounds and linear equalities they were.

The sensitivities behave the same way. A single shooting rollout compounds them across the
whole horizon, so the gradient of the cost with respect to an early input passes through
every subsequent step and grows or vanishes exponentially in unstable or stiff systems,
whereas multiple shooting confines each sensitivity to one interval. Multiple shooting
therefore buys conditioning and cheap state constraints with a larger and sparser problem,
and single shooting is preferable where the horizon is short enough for the compounding to
stay benign and the constraints bear on the input alone.

\begin{figure}[tb]
    \centering
    \includegraphics[width=0.9\linewidth]{figures/03_foundations/shooting}
    \caption{Direct single shooting against direct multiple shooting over one Control
    Horizon. Under single shooting the trajectory is one continuous rollout from the initial
    state and only the inputs are decision variables, so the dynamics hold at every iterate.
    Under multiple shooting the horizon is split into intervals, the state at the start of
    each interval is a decision variable of its own, and the gaps between the resulting
    trajectory pieces close only as the solver converges. The intermediate iterate drawn for
    the multiple shooting case is therefore not a valid trajectory of the system.}
    \label{fig:shooting}
\end{figure}

Whichever transcription is used, the result is an \ac{nlp} in finitely many variables,
solved by one of two families. \Ac{sqp} methods approximate the problem at the current
iterate by a \ac{qp} with a quadratic model of the Lagrangian and linearised constraints,
solve it for a step, and repeat, which suits receding horizon control because a warm start
from the previous solve often leaves few active-set changes to resolve. Interior-point
methods instead replace the inequality constraints by a barrier term and drive the barrier
parameter to zero, following a central path to the solution \cite{Waechter2006}. Interior-point methods scale
better with the number of constraints and are less able to exploit a warm start, since they
require a strictly feasible interior point and the previous solution typically sits on the
boundary \cite{Nocedal2006}.

\subsection{Output Feedback and Input-Output Models}
\label{sec:output_feedback}

The formulation \eqref{eq:ocp} assumes the initial state $\hat{\mathbf{x}}_k$ is known, and
in this application it is not. The plant carries $6N$ states plus the history of the delay
term, and the controller sees $M$ linear mixtures of a subset of them through
\eqref{eq:eeg}. The classical remedy places an observer in front of the controller,
estimates the state from the measurement sequence, and hands the estimate to the \ac{ocp} as
though it were exact. For linear systems the separation principle makes the arrangement
sound, but under nonlinear dynamics and active constraints no such principle holds, and
stability results for output-feedback \ac{nmpc} rest instead on the observer error decaying
fast enough relative to the controller \cite{Findeisen2003}. The construction also needs the
plant equations the observer is built from, which a data-driven setting does not provide.

The alternative taken here removes the estimation problem instead of solving it. An
input-output model of \ac{narx} type predicts the next measurement from a finite window of
past measurements and past inputs \cite{Chen1989}, so that
%
\begin{equation*}
    \mathbf{y}_{k+1} = \mathbf{f}_\mathrm{NARX}\!\left(\mathbf{y}_k, \dots,
    \mathbf{y}_{k-n_\mathrm{y}+1},\; \mathbf{u}_k, \dots, \mathbf{u}_{k-n_\mathrm{u}+1}\right),
\end{equation*}
%
with $\mathbf{y}_k$ the measurement at step $k$ and $n_\mathrm{y}$ and $n_\mathrm{u}$ the
numbers of past outputs and inputs retained. The window is the model state, $\mathbf{x}_k =
\left[\mathbf{y}_k, \dots, \mathbf{u}_{k-n_\mathrm{u}+1}\right]$, and it is assembled from
quantities the controller has recorded rather than inferred, so $\hat{\mathbf{x}}_k$ is
available exactly and neither an observer nor a separation argument is needed. Two
requirements replace the ones the observer imposed. The window has to be long enough to
carry the information about the plant state the model needs, which the delayed coupling of
\eqref{eq:coupling} makes concrete, and the model must be filled with a valid history before
its first prediction is meaningful. Chapter~\ref{ch:methodology} treats that Priming
requirement as part of the closed-loop schedule.

\subsection{Cost Function Design}
\label{sec:cost_design}

The stage cost decides what the controller considers success, and it separates into two
terms answering different questions,
%
\begin{equation*}
    \ell(\mathbf{x}, \mathbf{u}) = \ell_\mathrm{x}(\mathbf{x}) + \ell_\mathrm{u}(\mathbf{u}),
\end{equation*}
%
one scoring the predicted state or output and one scoring the input. The first expresses
what suppression means, the second what the stimulation is allowed to cost, and the two are
designed independently of each other.

For the state term the standard choice is tracking, a quadratic penalty on the distance
between the predicted quantity and a reference. Seizure suppression admits no single
reference trajectory, because any output waveform staying within the amplitude and spectral
envelope of healthy activity is acceptable and the many such waveforms are not ranked
against each other. A cost expressing a performance criterion directly rather than the
distance to a precomputed setpoint is called economic \cite{Ellis2014}, and a criterion
scoring how far a predicted quantity leaves a healthy envelope, while leaving it free
inside, is of that kind. Both forms are used in Chapter~\ref{ch:methodology}: quadratic
tracking terms drive predicted quantities towards zero wherever a meaningful target exists,
and economic terms take over wherever the objective is an envelope rather than a value.

The input term is the second half of the stage cost, and its choice does more than trade
effort against performance. Writing the input penalty as an $\ell_p$ norm of the input,
%
\begin{equation}
    \ell_\mathrm{u}(\mathbf{u}) = w \, \|\mathbf{u}\|_p^p,
    \label{eq:input_penalty}
\end{equation}
%
with $w > 0$ the weight, the quadratic case $p = 2$ penalises large currents quadratically
and small ones hardly at all, so its minimisers spread a small current over the whole
horizon and every electrode. The case $p = 1$ penalises with a constant slope down to zero
and drives components exactly to zero instead, the sparsity mechanism the lasso is built on
\cite{Tibshirani1996}. The limit of the sequence is the $\ell_0$ pseudo-norm counting the
nonzero entries, which is what a sparsity requirement actually means and which is neither
convex nor differentiable. The link between the two is that on a bounded set the $\ell_1$
norm is the convex envelope of $\ell_0$, the largest convex function bounding it from below
\cite{Boyd2004}, so $\ell_1$ is the tightest convex surrogate for sparsity available rather
than a merely convenient smooth stand-in. Under the sparsity motivation of
Section~\ref{sec:mi_mpc}, where the quantity to be economised is delivered charge, the
progression from $p = 2$ through $p = 1$ to the $\ell_0$ limit forms a single axis along
which the controller is asked for progressively more concentrated stimulation.

\subsection{Mixed-Integer MPC}
\label{sec:mi_mpc}

Taking the $\ell_0$ limit seriously changes the class of the optimisation problem. A
stimulation waveform delivering charge only in isolated bursts is described by a binary
decision per step, $\delta_{j|k} \in \{0, 1\}$, indicating whether a pulse is applied,
together with a continuous amplitude admissible only when $\delta_{j|k} = 1$. The coupling
between the two is the standard big-$M$ implication
%
\begin{equation*}
    -\delta_{j|k} \, u_\mathrm{max} \leq u_{p,j|k} \leq \delta_{j|k} \, u_\mathrm{max},
    \qquad p = 1, \dots, P,
\end{equation*}
%
with $u_{p,j|k}$ the predicted current at electrode $p$ and $u_\mathrm{max}$ the
per-electrode bound, which forces the amplitude to zero on steps the controller elects not
to stimulate and leaves it free otherwise. A budget of $n_\mathrm{p}$ pulses per horizon is
then the linear constraint $\sum_{j=0}^{H-1} \delta_{j|k} \leq n_\mathrm{p}$, and the
sparsity the $\ell_1$ penalty of \eqref{eq:input_penalty} only encourages becomes something
the problem enforces exactly. Systems mixing continuous dynamics with logical decisions of
this kind are the subject of the mixed logical dynamical framework \cite{Bemporad1999}.

Problem~\eqref{eq:ocp} extended this way is a \ac{miqp} when the underlying problem is a
\ac{qp} and a \ac{minlp} in the nonlinear case treated here. Both are solved by
branch-and-bound: a relaxation with the integrality dropped supplies a lower bound, a
fractional variable is selected and the problem split into the two subproblems fixing it to
zero and to one, and subtrees whose bound is worse than the best known integer solution are
pruned \cite{Belotti2013}. Pruning is what makes the method usable, but it guarantees
nothing, and the worst case remains the $2^H$ binary combinations over the horizon. The cost
is far more damaging in receding horizon control than in offline planning, because the
problem is solved not once but at every sampling instant, and because a \ac{minlp} node
relaxation is itself a nonconvex \ac{nlp} whose local minimiser is not a valid bound. Move
blocking, which constrains the input to be constant across groups of consecutive steps and
so reduces the number of binary variables below $H$, is the usual lever for keeping the tree
small enough \cite{Cagienard2007}.

The alternative to solving the integer problem is to relax it, solve the resulting
continuous problem, and recover a binary solution from the relaxed one. Relaxing
$\delta_{j|k} \in \{0,1\}$ to $\delta_{j|k} \in [0,1]$ returns the \ac{minlp} to an \ac{nlp}
of the kind Section~\ref{sec:transcription} already handles, at the price of a solution
physically meaningless where it is fractional, a pulse being either delivered or not. A
rounding scheme has to follow, and its value rests on whether the trajectory of the rounded
solution can be bounded against that of the relaxed one; bounds of this kind exist for
mixed-integer optimal control and tighten as the discretisation is refined \cite{Sager2012}.
Which relaxation and rounding to adopt is left open here, because the choice interacts with
the prediction model: a relaxation over pulse timing rather than pulse presence needs a
model evaluable between sampling instants, which a discrete-time predictor is not.
Chapter~\ref{ch:methodology} returns to the question once the prediction model is fixed.

\section{Data-Driven Dynamics Identification}
\label{sec:sysid}

\subsection{Multilayer Perceptrons (MLP)}


\section{Short-Time-Fourier Transform} % location of this section not sure


% the whole pipeline: windowing hop, kernel smooting, frequency resolution, etc.
\chapter{Methodology}
\label{ch:methodology}

% Chapter introduction, without a separate section:
% - Formulate the suppression problem, available measurements, admissible stimulation,
%   and the role of a learned Predictor in MPC.
% - Introduce the investigated choices: waveform versus spectral prediction, training
%   objectives, and continuous versus pulse-based stimulation.
% - Explain how these choices connect and introduce the chapter's order.
% Editorial boundary: Chapter 3 supplies general theory. This chapter formulates the
% methods and motivates design choices. Chapter 5 tests them and reports the evidence.
% Introduce tunable parameters and selection criteria here; put search ranges, budgets,
% results, and selected configurations in Chapter 5. Routine tuning can go in an appendix.

\section{Closed-Loop Neurostimulation Architecture}

% i don't like this whole first paragraph. Feels too much like AI
Chapter~\ref{ch:foundations} introduced the plant, electrical interfaces, prediction
models, and receding-horizon control separately. This section connects them into the
closed-loop setting used in this thesis. The setting follows the EEG-feedback structure
of Yu et al.~\cite{Yu2024}, but replaces their amplitude trigger with a prediction model % we don't need to mention what controller they used. This should be in the introduction not here.
and \ac{mpc}. Figure~\ref{fig:closed_loop_architecture} shows the resulting signal path.
Only scalp \ac{eeg} and previously applied electrode currents enter the controller.
The regional activity retained to assess seizure propagation is not available for feedback.

At a controller update, the signal-conditioning stage passes the latest causal % its not "causal signal conditioning" its a low pass filter.
measurement to the prediction model. The \ac{mpc} problem evaluates candidate current
sequences and applies the first current vector of its solution. A zero-order hold keeps
this vector constant until the next update. The stimulation projection converts the
electrode currents into regional membrane perturbations, after which the plant evolves
and produces the next measurement. Waveform and spectral controllers use this same
physical loop. Their update rates differ because a waveform sample is available after
each decimation interval, whereas a spectral quantity is available only after a new
analysis frame has been formed.

\begin{figure}[htbp]
    \centering
    \begin{tikzpicture}[
        block/.style={
            draw,
            rounded corners=2pt,
            align=center,
            minimum height=1.3cm,
            text width=1.9cm,
            font=\small
        },
        signal/.style={->, thick},
        evaluation/.style={->, dashed, thick}
    ]
        \node[block] (plant) at (0,4.2) {Whole-brain seizure plant};
        \node[block] (eeg) at (2.8,4.2) {62-channel EEG projection};
        \node[block] (conditioning) at (5.6,4.2) {Causal signal conditioning};
        \node[block] (control) at (8.4,4.2) {Prediction model and MPC};
        \node[block] (field) at (8.4,2.0) {ROAST field projection};
        \node[block] (assessment) at (2.8,0) {Seizure-spread assessment};

        \draw[signal] (plant) -- node[above, font=\scriptsize] {$\mathbf{y}_\mathrm{LFP}$} (eeg);
        \draw[signal] (eeg) -- node[above, font=\scriptsize] {$\mathbf{y}_\mathrm{EEG}$} (conditioning);
        \draw[signal] (conditioning) -- node[above, font=\scriptsize] {$\tilde{\mathbf{y}}_k$} (control); % is tilde what we use for denoting estimates/filtered values?
        \draw[signal] (control) -- node[right, font=\scriptsize] {$\mathbf{u}_k$} (field);
        \draw[signal] (field) -- (0,2.0) -- node[left, font=\scriptsize] {$\mathbf{s}_k$} (plant);
        \draw[evaluation] (plant) -- (0,0) -- node[above, font=\scriptsize] {evaluation only} (assessment);
    \end{tikzpicture}
    \caption{Closed-loop neurostimulation architecture. Solid arrows denote the causal
        feedback loop. The dashed branch carries regional activity used only for post hoc
        evaluation and does not provide the controller with source-level information.}
    \label{fig:closed_loop_architecture}  %  I don't like this picture. It should show the blocks in a standard closed-loop control block diagram layout. With the EEG projection and ROAST field projection just being multiplication blocks and holding their matrix symbols. There needs to be a "reference"/goal. The seizure spread assessment is not a part of it at all
\end{figure}

\subsection{Plant Configuration and Seizure Dynamics}

The loop in Figure~\ref{fig:closed_loop_architecture} is closed around the biophysical
model of Section~\ref{sec:excitability}. Its plant contains 76 Jansen--Rit regions coupled
by the structural weights and tract lengths of the connectome. Each region contributes  % what connectome? mention that it is the TVB ...
six states, so the instantaneous neural-mass state has dimension 456. Fibre lengths are  % the neural-mass state is irrelevant information. This should just be a description of our setup and what decision we took.
converted into delays using the uniform conduction speed
$c=\SI{50}{\milli\metre\per\milli\second}$, and each delay is rounded to the nearest plant
integration step.

The regional seizure configuration adopts the temporal-lobe partition of Yu et.\.al.~\cite{Yu2024}. % and why is that partition important?
The left hippocampus, left parahippocampal cortex, and left amygdala form  % we need to explain why these regions were chosen.
the \ac{ez}, with $A_i=\SI{3.6}{\milli\volt}$. The left inferior temporal cortex and left
ventral temporal cortex form the initial \ac{pz}, with
$A_i=\SI{3.4}{\milli\volt}$. All remaining regions use the background value
$A_i=\SI{3.25}{\milli\volt}$. The global coupling factor $K=0.60$ and noise intensity
$\sigma=\SI{280}{\per\second}$ were calibrated jointly because both parameters affect the
time at which activity escapes the \ac{ez} and recruits the \ac{pz}. The calibration
criterion rewards early recruitment within the \ac{ez}, later recruitment of the \ac{pz}
and left hemisphere, and little recruitment of the right hemisphere. The calibration
results are reserved for the appendix.

Each simulation starts from the settled healthy state and advances the stochastic delayed
dynamics with the Heun method at $\Delta t=\SI{0.1}{\milli\second}$. One Gaussian increment
is used in both stages of a step. For a standard normal random vector
$\boldsymbol{\xi}_k$, the stochastic increment on the excitatory loop is therefore
$\sigma\sqrt{\Delta t}\,\boldsymbol{\xi}_k$. This fine grid resolves the neural-mass
dynamics and conduction delays. It is not the sampling grid presented to the controller.

\subsection{EEG Measurement and Stimulation Geometry}

The plant configuration fixes the internal seizure dynamics, while closing the loop
requires separate maps for measurement and stimulation. For measurement, the surface
projection supplied with The Virtual Brain maps 16,384 cortical vertices to 62 scalp
electrodes~\cite{SanzLeon2015}. Summing the columns belonging to each anatomical region
reduces this operator from $62\times16{,}384$ to the regional lead field
$\mathbf{L}_\mathrm{EEG}\in\mathbb{R}^{62\times76}$. Equation~\eqref{eq:eeg} then maps the
76 regional \acp{lfp} to the measured scalp signals. No additional measurement noise is
added. Stochastic variability enters through the plant noise described above.

The regional neural-mass output is only proportional to a dipole moment, and neither this
proportionality nor the lead-field data establish an absolute voltage scale. Consequently,
the simulated scalp \ac{eeg} is stated in arbitrary amplitude units. Thresholds and healthy
references are constructed from signals produced by the same forward model, so their
relative scale remains consistent.

The stimulation interface uses the TP9--CP5--Ex8 montage. It extends the CP5--Ex8
placement used by Yu et al.~\cite{Yu2024} with TP9 as a second controllable scalp site.
Ex8 is the return position below the inion defined in the 10--05 cap. During field
generation, ROAST substituted its TPP9 position for the requested TP9 label because this
was the nearest available cap position. The thesis retains TP9 as the requested montage
label and records the substitution here.

ROAST computes one electric-field basis solution for a \SI{1}{\milli\ampere} current at
each controllable scalp electrode against Ex8 on the MNI152 head model~\cite{Huang2019}.
Linear superposition gives the field for any balanced current vector. Let
$\mathbf{E}_{p,i}$ denote the field at the centroid of region $i$ due to the basis current
at electrode $p$, and let $\mathbf{n}_i$ be the outward cortical normal assigned to that
region. The regional membrane perturbation is
%
\begin{equation}
    s_i(t)
    = \lambda_\mathrm{p}\sum_{p=1}^{3}u_p(t)\,
      \mathbf{E}_{p,i}^{\mathsf{T}}\mathbf{n}_i,
    \qquad
    \lambda_\mathrm{p}=\SI{0.35}{\milli\metre} \, .
    \label{eq:regional_stimulation_projection}
\end{equation}
%
The current $u_p$ belongs to electrode $p$, and $\lambda_\mathrm{p}$ is the assumed somatic
polarisation length. The projection in~\eqref{eq:regional_stimulation_projection} uses a
value within the range measured for deep-layer
pyramidal cells~\cite{Radman2009}. Since
$\mathbf{E}_{p,i}^{\mathsf{T}}\mathbf{n}_i$ has units of volts per metre, equivalently
millivolts per millimetre, $s_i$ has units of millivolts. It enters the pyramidal sigmoid as
specified in Section~\ref{sec:electro}. The controller imposes
$|u_p|\leq\SI{2}{\milli\ampere}$ and $\sum_{p=1}^{3}u_p=0$. No slew-rate constraint is used.

The transpose of $\mathbf{L}_\mathrm{EEG}$ cannot supply this stimulation map. Reciprocity
applies when measurement and stimulation use the same electrodes, source locations, and
orientations~\cite{Malmivuo1995}. Ex8 is absent from the 62-channel recording montage, and
the ROAST map samples a volumetric electric field along regional cortical normals. These
conditions differ from those used to construct the \ac{eeg} forward operator.

\subsection{Filtering and Sampling}

The spatial projections determine which signals the controller can measure and apply;
filtering determines which temporal content reaches it. In the simulated seizure regime,
the dominant oscillation lies near \SIrange{3}{4}{\hertz}, and most seizure-related power
lies below \SI{12}{\hertz}, as illustrated in Figure~\ref{fig:bifurcation}. The controller input is
therefore sampled at \SI{50}{\hertz}, retaining frequencies up to the
\SI{25}{\hertz} Nyquist limit while reducing the number of prediction steps.

Before decimation from the \SI{10}{\kilo\hertz} plant grid, every \ac{eeg} channel passes
through a causal fourth-order Butterworth low-pass filter with a \SI{25}{\hertz} cutoff.
A causal filter is required because future samples are unavailable during feedback; a
zero-phase offline filter would give the controller non-causal information. The filtered
signal is then decimated by 200, producing one waveform sample every
\SI{20}{\milli\second}. The same filtering and decimation are applied to identification
data and during closed-loop operation.

The \SI{25}{\hertz} cutoff defines the anti-aliasing stage, not the frequency range used by
a later spectral cost or prediction representation. A waveform controller may update on
the \SI{20}{\milli\second} sample grid. A spectral controller instead updates when its next
causal analysis frame becomes available, so its interval also depends on the frame hop.
Section~\ref{sec:prediction_representations} next defines these two prediction
representations and their time grids.

\section{Prediction Representations and Models}
\label{sec:prediction_representations}
% Motivate the central question: what should be predicted to support seizure suppression?
% Separate the output representation from the model class used to predict its dynamics.

\subsection{Waveform Prediction}
% Define the future EEG waveform to be predicted, the measured history, and the past
% and candidate future currents conditioning the prediction. Motivate this representation.

\subsection{Spectral Observable Prediction}
% Define the spectral Observable and why it is a candidate alternative to waveform prediction. The reason is: if we are only predicting
% Refer to Chapter 3 for STFT theory; specify the representation used in this method.
% Introduce Segment length, Frame hop, frequency selection, pooling, and normalization.
% Explain which information each geometry choice retains or discards and its time scale.
% Keep geometry parameterized here; its empirical selection belongs in Chapter 5.
% QUESTION: Is geometry selected for prediction accuracy, closed-loop suppression,
% temporal responsiveness, or an explicit tradeoff? Define "optimal" relative to that goal.

\subsection{Predictor Model Classes}
% Define the shared autoregressive input-output formulation and multi-step Rollout.
% Introduce the nonlinear MLP and linear/DMD variants and their modeling assumptions.
% Relate DMD to the Koopman perspective only as far as needed to motivate the chosen model.
% Explain which model classes are studied for each representation and why.
% Describe initialization from measured history only to establish a valid prediction.

\section{Predictor Training and Selection}
% Explain how the prediction tasks become identification problems and how candidates
% are selected. This section defines procedures; Chapter 5 reports their outcomes.

\subsection{Identification Data and Prediction Task}
% Describe simulated training data, seizure regimes, and stimulation used to identify
% the response to currents rather than only autonomous dynamics.
% Define history/target construction, prediction duration, and normalization.
% Explain separation of training, selection, and final evaluation data at trajectory level.
% Put experimental dataset sizes and scenario-specific settings in Chapter 5.

\subsection{Waveform and Spectral Training Losses}
% Formulate waveform error, spectral error, and the combinations to be investigated.
% Distinguish a spectral Loss on predicted waveforms from prediction of spectral Observables.
% For the former, STFT geometry defines error scoring; for the latter, it defines the target.
% Introduce the Loss geometry, weights, and scored Spans without claiming selected values.
% QUESTION: Should these two STFT geometries be selected independently, and what must
% remain comparable when evaluating the resulting Predictors?

\subsection{Training Strategies and Selection Criteria}
% Describe fitting procedures and curriculum strategies that are part of the investigation.
% Define the criterion used to select geometry, Loss settings, and model configurations.
% QUESTION: Which choices are scientific comparisons and which are routine tuning?
% QUESTION: Is selection based on open-loop prediction or closed-loop control performance?
% Chapter 5: give search spaces/budgets, compare outcomes, and justify the configurations
% carried into later experiments. Reserve independent evaluation data for final comparisons.

\section{MPC Objectives and Continuous Stimulation}
% Turn the trained Predictor into a controller by defining the optimization problem.

\subsection{Decision Variables and Control Constraints}
% Define the Control Horizon, candidate electrode currents, and predicted dynamics.
% Specify Kirchhoff's law and the Control Budget imposed in the study.
% State the current parameterization and how prediction and control time grids relate.

\subsection{Seizure-Suppression Costs}
% Formulate waveform-based and spectral suppression Costs and their intended behavior.
% Explain construction of the healthy spectral reference and normalization of Cost terms.
% Keep the controller's suppression Cost distinct from the Predictor's training Loss
% and from the post hoc Closed-Loop Evaluation Score used in Chapter 5.

\subsection{Stimulation Penalties}
% Define L2 and L1 current penalties, weight scaling, and the behavior to be compared.
% Describe the smooth L1 surrogate where it is used in the evaluated method.
% End with the complete continuous optimization problem, collecting dynamics, Costs,
% and constraints before extending the formulation to pulses.

\subsection{Numerical Solution and Receding-Horizon Application}
% Explain the transcriptions and solution methods studied, including single/multiple
% shooting and NARX/PCMS formulations if retained as part of the investigation.
% State how Costs spanning multiple samples or Frames enter the optimization.
% Avoid a software architecture account.
% Chapter 5 contains solver comparison protocols and results, not first definitions.

\section{Pulse-Based Control}
% Formulate pulse-based stimulation as an alternative to continuous stimulation.
% Reuse the continuous problem's Predictor, suppression Cost, and current constraints.

\subsection{Pulse Definition and Scheduling Variables}
% Define a pulse before introducing its decision variables and admissible schedules.
% QUESTION: Is the objective to restrict current magnitude, total stimulation time,
% the number of distinct pulses, or some stated combination?
% QUESTION: Does a binary variable denote an active interval or a pulse onset?
% Summing active steps counts enabled intervals; consecutive enabled steps can form one pulse.
% QUESTION: Are pulse shape, duration, polarity, and electrode pattern fixed?
% Which of amplitude, duration/duty cycle, and timing should be optimized?

\subsection{Mixed-Integer Formulation}
% Couple binary scheduling decisions to admissible currents and define the pulse budget.
% State which decisions are shared across electrodes and which are electrode-specific.
% Derive the resulting optimization class from the chosen Predictor, Cost, and constraints.
% Explain the solution approach and any schedule simplifications needed for tractability.
% QUESTION: What is the smallest pulse formulation that tests the intended idea before
% adding amplitude or duty-cycle optimization?

\subsection{Continuous Relaxation and Schedule Recovery}
% Define a relaxation of binary activation variables and how a deliverable schedule
% would be recovered. Explain how feasibility and predicted suppression would be rechecked.
% QUESTION: Which recovery rule preserves the chosen pulse budget and current constraints?
% Distinguish this procedure from continuous optimization of pulse onset times.

\subsection{Continuous Pulse-Timing Alternative}
% Formulate onset times as continuous decisions (if this alternative remains in scope)
% Explain how candidate pulses occurring within a sampling interval affect prediction.
% Develope a continuous-time Predictor required
% Compare the formulations conceptually with L2/L1 continuous stimulation here.
% Chapter 5 compares their suppression, stimulation use, and computational cost.

\chapter{Experiments and Evaluation}
\label{ch:experiments}

% Generally: Each experiment should have its own question, setup, results, and interpretation.
% what follows is just an idea. Don't take too rigidly
% need to prove how purely waveform prediction - cost and cost fails
% --------


% Chapter introduction:
% - Explain how the experiments answer the design questions formulated in Chapter 4.
% - Establish the sequence: select geometry and training settings, compare Predictors
%   and numerical methods, then evaluate closed-loop stimulation.
% - End each study with the evidence supporting the choices used in subsequent studies.
% Editorial boundary: refer to Chapter 4 for method definitions. Give experimental
% settings, search spaces/budgets, results, and interpretation here.
% For each comparison, state the question, what varies, what is held fixed, and the
% criterion for judging the outcome. Do not assume that the expected improvement occurs.

\section{Experimental Setup and Evaluation Criteria}

\subsection{Simulation Scenarios and Data Splits}
% Specify Plant configurations, simulation durations, stimulation scenarios, dataset
% sizes, and independent repetitions. Refer to Chapter 4 for construction of the Plant.
% Describe healthy and unstimulated seizure activity as reference conditions.
% Include evidence for Plant calibration here if that calibration is an empirical study.
% State which trajectories/seeds are used for training, selection, and final evaluation.
% Use matched conditions across methods where possible and report variability across runs.
% QUESTION: Which seizure regimes and stimulation conditions must the conclusions cover?

\subsection{Prediction Accuracy and Seizure Metrics}
% Define evaluation measures and aggregation over channels, time, and repetitions.
% Include waveform and spectral prediction errors, Seizure State, and Seizure Burden
% as relevant. Distinguish sensor-space scores from regional ground truth.
% Assess prediction error over physical forecast duration, not just the number of steps.
% Keep training Loss, controller Cost, and Closed-Loop Evaluation Score distinct.
% QUESTION: What operational criterion defines a useful prediction duration for control?

\subsection{Stimulation Use and Computational Performance}
% Define measures for current magnitude, stimulation duration, distinct pulse count,
% and constraint satisfaction according to the formulations in Chapter 4.
% State the computational setup and how solve time, failures, and solution quality are measured.
% QUESTION: Is a controller update deadline part of the study, or is runtime a comparison
% measure only? Do not claim real-time feasibility without testing a stated deadline.

\section{Selection of Spectral Observable Geometry}
% QUESTION: Which STFT geometry makes the spectral Observable useful for prediction
% and control, under the selection criterion established in Chapter 4?

\subsection{Geometry Search and Comparison Protocol}
% Give candidate Segment lengths, Frame hops, and other geometry parameters varied.
% State the Predictor class, fitting procedure, search budget, and selection data.
% Account for changes in target dimension, Frame rate, and physical prediction duration.
% QUESTION: What common evaluation reference makes scores comparable when changing
% the geometry changes the target itself? Avoid ranking solely by native training Loss.

\subsection{Prediction Accuracy and Temporal Responsiveness}
% Compare prediction quality and preservation of seizure onset/propagation dynamics.
% Report tradeoffs associated with temporal and spectral resolution and smoothing.
% If selection uses closed-loop performance, include that evidence on selection scenarios
% here and reserve independent scenarios for the final closed-loop comparison.
% QUESTION: Does a geometry that is easier to predict still resolve changes needed for control?

\subsection{Selected Geometry}
% Justify the configuration carried into later experiments using the declared criterion.
% Report sensitivity to nearby choices and variability rather than only the best trial.
% State what the evidence establishes about this choice and where it remains uncertain.

\section{Waveform Predictor Training-Loss Comparison}
% QUESTION: How do waveform Loss, spectral Loss, and curriculum training affect the
% accuracy and useful duration of waveform predictions?

\subsection{Loss Geometry and Training Configurations}
% Compare the Loss combinations and curricula defined in Chapter 4, including the
% intended comparison of STFT training with curriculum plus STFT training.
% Report searches over spectral Loss geometry, weights, and Spans where investigated.
% This geometry scores waveform predictions; it need not equal the spectral Observable geometry.
% Control model capacity, data, training effort, and repetitions, or explain differences.
% QUESTION: Which comparisons isolate the effect of the Loss, geometry, and curriculum
% rather than changing all three at once?

\subsection{Prediction Results and Training Strategy Selection}
% Evaluate waveform and spectral errors on common criteria across training objectives.
% Examine multi-step error and retention of seizure dynamics and stimulation responses.
% Justify the Loss geometry and training strategy used in the Predictor comparison.
% Assess their consequences for control in the closed-loop study below; if control is
% the selection criterion, report selection evidence here separately from final evaluation.

\section{Predictor Comparison}
% QUESTION: What is gained or lost by predicting waveforms versus spectral Observables,
% and by using linear/DMD versus nonlinear MLP dynamics?

\subsection{Comparison Protocol and Model Selection}
% List the representation/model combinations studied and the configurations selected above.
% Report remaining hyperparameter searches and comparable selection budgets.
% Separate the representation comparison from the model-class comparison where possible.
% QUESTION: Are the selected geometry and Loss settings transferable across model classes,
% or does each class require its own selection? State the fairness criterion used.

\subsection{Prediction Quality over the Control Horizon}
% Compare errors at matched physical forecast durations and across relevant seizure regimes.
% Use shared spectral evaluation criteria where waveform errors cannot be compared directly.
% Assess response to candidate stimulation, as well as unstimulated seizure evolution.
% Identify Predictors carried into closed-loop evaluation without equating prediction
% accuracy with demonstrated control performance.

\section{Numerical Solution Comparison}
% QUESTION: Which formulation and solver solve the control problems reliably at the
% Predictor dimensions and Control Horizons used in the study?
% Keep this section proportional to its scientific role; shorten it if solver selection
% is only a supporting choice. Definitions belong in Chapter 4.

\subsection{Single Shooting, Multiple Shooting, and PCMS}
% Compare the selected transcriptions, including the NARX and PCMS formulations.
% Hold the underlying Predictor, Cost, constraints, and initial control problem fixed.
% State initialization, termination tolerances, and treatment of whole-horizon Costs.

\subsection{Solution Quality, Reliability, and Runtime}
% Compare attained Cost, constraint violations, failure frequency, and solve time.
% Examine the effect of Predictor class and Control Horizon where relevant.
% Justify the numerical methods used in the subsequent closed-loop experiments.
% QUESTION: Does a faster solve preserve the objective and feasibility of the original problem?

\section{Closed-Loop Control Evaluation}
% Test whether the choices above lead to seizure suppression under admissible stimulation.
% Use final evaluation scenarios that were not used to select configurations.

\subsection{Evaluation Protocol and Reference Conditions}
% Specify controller configurations, Control Horizons, update intervals, and evaluation durations.
% Include the unstimulated reference and justify any additional control baselines.
% State how controller Cost weights and budgets were selected before final evaluation.
% QUESTION: Which comparisons isolate the Predictor, suppression Cost, and stimulation
% formulation, and which compare the best selected combination for each method?

\subsection{Prediction Representation and Suppression Objective}
% QUESTION: Do the open-loop prediction advantages translate into lower Seizure Burden?
% Compare waveform and spectral Predictors and the suppression Costs defined in Chapter 4.
% Relate sensor-space improvement to regional seizure outcomes and stimulation use.
% Include focused training-Loss comparisons if needed to establish their effect on control.

\subsection{Continuous and Pulse-Based Stimulation}
% QUESTION: How do L2/L1 continuous control and pulse-based formulations trade seizure
% suppression against current magnitude, stimulation time, pulse count, and computation?
% Compare the mixed-integer method, relaxation with schedule recovery, and continuous
% pulse-timing method included in the thesis, under explicitly comparable restrictions.
% Evaluate recovered schedules as delivered to the Plant, not just their relaxed objectives.
% QUESTION: What changes after schedule recovery in feasibility and achieved suppression?
% QUESTION: Does optimizing amplitude or duty cycle add benefit over the simpler pulse formulation?

\subsection{Tradeoffs and Limits of the Findings}
% Bring together suppression, stimulation use, and computational reliability across runs.
% Identify where prediction-based selection agrees with closed-loop performance and where it fails.
% Distinguish conclusions supported across scenarios from those specific to the simulated setup.
% Close by answering the chapter's questions and preparing the thesis conclusion.

\chapter{Conclusion and Discussion}
\section{Summary of Key Findings}
\section{Limitations and Critical Discussion}
\section{Future Directions}


\cite{Duden}

\appendix

\cleardoublepage
\phantomsection
\addcontentsline{toc}{chapter}{\bibname}

\printbibliography

\end{document}
